\section{Data}\label{sec:data}

UniPlan uses one real data source and several synthetic ones (Table~\ref{tab:data}). The real source is the public University of Sydney timetable for Semester~2, 2026~\cite{usyd_timetable}. A small scraper fetched only the units we needed on 24--25 September 2026, at most one request a second, and kept no staff names. The snapshot covers \dUnits{} units, \dMeetings{} dated class meetings, \dOptions{} class options and \dGroups{} activity groups. Classes were expanded to their real teaching dates, so alternate weeks and one-off assessments are handled. Seat numbers are as of the snapshot, taken mid-semester, so classes that were full then are excluded from every count and test (for ENGG2112 and AMME2301, \ttTwoFull{} of \ttTwoOptions{} classes). Walking times use the \dBuildings{} buildings' positions on the public campus map.

\begin{table}[h]
	\caption{Datasets. Everything below the line was written or labelled by AI models; no real student wrote a sentence, answered a question or labelled a pair.}
	\label{tab:data}
	\centering\small
	% Generated by report/scripts/make_numbers.py from the raw result files. Do not edit by hand.
{\small
\begin{tabular}{@{}>{\raggedright\arraybackslash}p{1.45in}>{\raggedright\arraybackslash}p{1.6in}>{\raggedright\arraybackslash}p{1.7in}>{\raggedright\arraybackslash}p{1.0in}@{}}
\toprule
Dataset & Size & Written or labelled by & Used for \\
\midrule
\multicolumn{4}{@{}l}{\textit{Real}} \\
Timetable snapshot, 2026 S2 (24 and 25 Sep 2026) & 7 units, 1{,}240 meetings, 242 classes & University of Sydney & Scheduling \\
Building locations & 19 buildings & University campus map & Walking time \\
\midrule
\multicolumn{4}{@{}l}{\textit{Synthetic (AI-written)}} \\
Generated clauses & 1{,}781 & Qwen2.5-7B; sample graded by Claude & Parser training \\
Hard clauses & 141 + 201 (two rounds) & Claude & Parser training \\
Parsing tests v1, v2 & 142 + 126 sentences & Claude & Development \\
Parsing exam & 304 sentences, 461 clauses & Two fresh Claude agents & Final test \\
Personas & 140 (120 train, 20 held out) & Claude & Ranker \\
Teacher labels & 2{,}471 pairs: 2{,}061 asked (1{,}853 consistent), 410 by code & Qwen2.5-14B or code; audited by Claude & Ranker training (deployed model: first 60 personas) \\
Ranking exam & 20 personas, 540 questions & Fresh Claude agents & Final test \\
\bottomrule
\end{tabular}}

\end{table}

\textbf{All evaluation data is synthetic.} The sentences, personas, preference labels and both final exams were produced by AI models, not by students. We kept the generators apart from the models under test where we could: Qwen2.5-7B wrote most of the parser's training clauses (Claude added \sClaudeTrain{} harder ones), while fresh Claude agents that saw no training data wrote every test set and exam; Qwen2.5-14B labelled the ranker's training pairs, while separate Claude agents answered the ranking exam. Each model was frozen before its exam was committed with a SHA-256 hash, and each exam was written by a new agent that saw only the labelling rules and was scored once (Appendix~\ref{app:provenance}). This independence is between model families, not between people (Section~\ref{sec:limits}).
